% Further consequences; no priority claim. Owned by /root/frobenius.
\paragraph{C.4r (Proved: a single integral numerical surface class detects the remaining cases).}
Let $A$ be an abelian variety of dimension $g\in\{2,3\}$ over an
algebraically closed field $k$ of arbitrary characteristic.
Put $X=A\times\widehat A$, write $P=c_1(\mathcal P_A)$, and set
\[
 N^r(X)=\mathrm{CH}^r(X)/{\equiv_{\mathrm{num}}},
 \qquad r=2g-2,
\]
with integral coefficients. We regard this torsion-free group as
a subgroup of $N^r(X)_{\mathbb Q}$.
The following are equivalent:
\begin{enumerate}
\item $A$ is an actual direct summand of a product of Jacobians;
\item the rational numerical class
\[
                 \gamma_A=\frac{P^{2g-2}}{(2g-2)!}
                 \in N^{2g-2}(X)_{\mathbb Q}
\]
belongs to the integral numerical lattice $N^{2g-2}(X)$.
\end{enumerate}
Thus we require an integral dimension-two cycle numerically
equivalent to $P^2/2$ for a surface, or to $P^4/24$ for a
threefold. Over $\mathbb C$ the same conditions are also
equivalent to integral algebraicity of the corresponding
integral Hodge class.

\emph{Proof.}
We first justify the integral lattice used in (2).
The rational numerical group is finite-dimensional: numerical
equivalence is coarser than homological equivalence for any
Weil cohomology. Choose finitely many integral complementary
cycle classes $b_1,\ldots,b_t$ whose intersection functionals
span its rational dual. Then
\[
 N^r(X)\longrightarrow\mathbb Z^t,\qquad
 z\longmapsto\bigl(\deg(zb_j)\bigr)_j
\]
is injective. Consequently $N^r(X)$ is a finite free abelian
group.

Suppose (1) holds. By \citet[Theorem 1.2 and the paragraph
immediately following it]{cyScaviaSuzuki2025}, every component
of $\operatorname{ch}(\mathcal P_A)$ is algebraic with
$\mathbb Z_\ell$ coefficients for every prime $\ell$; no
principal polarization on $A$ is required for this implication.
In characteristic $p>0$ this includes the crystalline cycle
map with source $\mathrm{CH}^r(X)\otimes\mathbb Z_p$, as
specified in their equation (1.2).
Map each such representative to $N^r(X)\otimes\mathbb Z_\ell$.
Compatibility of the cycle map with intersections shows that
its pairing with every $b_j$ agrees with that of $\gamma_A$.
The injectivity of the displayed pairing after extending
scalars to $\mathbb Q_\ell$ therefore gives
\[
 \gamma_A\in N^r(X)\otimes\mathbb Z_\ell
 \qquad\text{for every prime }\ell.
\]
For $\ell=p$ the same comparison takes place in the fraction
field of $W(k)$, with the intersection values lying in
$\mathbb Q_p$; the source already has $\mathbb Z_p$
coefficients. Intersecting these local lattices inside
$N^r(X)_{\mathbb Q}$ gives
\[
 N^r(X)_{\mathbb Q}\cap
       \bigcap_{\ell}\bigl(N^r(X)\otimes\mathbb Z_\ell\bigr)
       =N^r(X),
\]
which proves (2). This passage uses neither the Tate conjecture
nor equality of numerical and homological equivalence.

Conversely, choose an integral cycle representing the numerical
class in (2), and intersect it with the integral divisor class
$P$. The resulting integral curve cycle has numerical class
\[
 \frac{P^{2g-1}}{(2g-2)!}
       =(2g-1)\frac{P^{2g-1}}{(2g-1)!}.
\]
Apply the integral Matsusaka factorization
\eqref{eq:c4-matsusaka-integer} on $X$ with the unimodular
Poincar\'e line bundle. Since its Euler characteristic is
$(-1)^g$, we obtain $(2g-1)\in I_J(X)$, after adjusting the
sign of one homomorphism, and therefore $(2g-1)\in I_J(A)$.
For $g=2$, C.4n already gives $2\in I_J(A)$, so
$1=3-2\in I_J(A)$. For $g=3$, C.4p gives
$6\in I_J(A)$, so $1=6-5\in I_J(A)$.
The explicit splitting of C.4h proves (1).
The Matsusaka homomorphism used here depends only on the
numerical curve class, as in C.4i: its Fourier correspondence
is a divisor class, and numerical triviality for divisors
annihilates their homomorphism to the dual abelian variety.

Finally, over $\mathbb C$, (1) implies integral algebraicity
of every component of $\operatorname{ch}(\mathcal P_A)$ by
\citet[Theorem 1.4]{cyVoisin2024} and
\citet[Theorem 1.1]{cyBeckmannFortman2023}. Conversely an
integral cohomological representative supplies the numerical
representative in (2). These divided powers are integral in
cohomology; this does not by itself divide the indicated
intersection cycle in the integral Chow group.
\hfill$\square$

The numerical integrality assertion in (2) remains to be proved
or disproved in the residual cases. Rational algebraicity of
the cohomological Fourier isometry of an abelian surface alone
does not supply the integral numerical cycle required here.

\paragraph{C.4s (Known consequence: every supersingular surface is affirmative).}
Every supersingular abelian surface over an algebraically closed
field of positive characteristic is an actual direct summand of
a product of Jacobians, including in characteristic two.
Indeed, \citet[Lemma 6.2, p.~185]{frC4FuLi2021} prove that it
admits a principal polarization. Their Proposition 6.1,
p.~184, explicitly includes characteristic two in the
N\'eron--Severi lattice calculation used by the proof.
The principally polarized surface case C.4f then applies.
This does not settle arbitrary nonsupersingular surfaces whose
polarization Euler characteristics are all even.
